% ============================================================
%  7. Ilustrativni prototip
%  Budžet: ~2 str. Pogl. 7 i 8 zajedno ≈ 12% opsega.
%  Struktura zaključana 2026-08-13.
%
%  ⚠ KOSTUR PREPISAN 2026-08-13, tekst napisan 2026-08-13 iz izvedenog
%    prototipa u prototip/. Raniji kostur je bio ZASTARJEO (ljestvica
%    brancheva, harness, tablice iz mjerenja). Ne vraćati.
%
%  ULOGA POGLAVLJA, JEDNA REČENICA: prototip postoji da bi ćelije matrice
%  u pogl. 8 iz razine (c) prosudba prešle u razinu (b) pokazano.
%  NE služi za mjerenje. NE dokazuje izvedivost proizvoda.
%
%  ⚠ SVI ISPISI U OVOM POGLAVLJU SU STVARNI, prepisani iz izlaza naredbi
%    `php artisan scenario` i `php artisan state` od 2026-08-13.
%    Ako se prototip mijenja, ISPISI SE MORAJU PONOVO UZETI.
%    Sat prototipa je fiksiran na 2026-09-01 09:00:00 pa su ponovljivi.
% ============================================================
\chapter{Ilustrativni prototip}
\label{ch:implementacija}

Obrasci iz odjeljka~\ref{sec:dizajn} opisani su riječima. Dio tvrdnji o njima
provjerljiv je jedino tako da se obrazac izvede i da se vidi u kojem stanju
sustav završi. Prototip služi tome, i ničemu drugom. Ne mjeri, ne dokazuje da
je proizvod izvediv i ne pokazuje sukladnost s propisom. Njegov jedini posao je
da tvrdnja koja bi inače ostala prosudba autora postane tvrdnja koju čitatelj
može pokrenuti. Poglavlje~\ref{ch:evaluacija} tu razliku vodi izrijekom.

% ------------------------------------------------------------
\section{Opseg i granice}
\label{sec:opseg-prototipa}
% ⚠ TVRDA GRANICA (2026-08-12).
%   Ovo se piše PRVO, prije opisa, da čitatelj zna što ne traži dalje.
%   Stack: PHP 8.5.8, Laravel 13.25.0, SQLite. Provjereno 2026-08-13.

Prototip pokriva jednu granicu, onu između poslovnog sustava i njegove
pristupne točke. To je granica C1 prema C2 iz odjeljka~\ref{sec:arhitektura},
ista koja u odjeljku~\ref{sec:failure-modes} nosi slučaj E1. Nijedan drugi dio
tijeka nije izveden.

Popis onoga što prototip nema je dulji od popisa onoga što ima, i tako je
namjerno:

\begin{itemize}[nosep]
  \item bez sastavljanja i provjere UBL dokumenta,
  \item bez Schematron pravila hrvatskog suženja norme,
  \item bez profila AS4 i bez ikakvog stvarnog prijenosa,
  \item bez elektroničkog potpisa i bez infrastrukture javnog ključa,
  \item bez otkrivanja adrese primatelja,
  \item bez mjerenja kašnjenja i bez ubrizgavanja kvarova u više servisa.
\end{itemize}

Pristupna točka nije izvedena nego zamijenjena \strani{krnjim
dvojnikom}{stub}. On na poziv vraća jedan od tri ishoda: potvrdu, izostanak
odgovora ili šifru S008. Prototip koji bi provjeravao UBL i govorio AS4
pokazivao bi sukladnost, a predmet je rukovanje greškom. Za njega je dovoljna
jedna granica i tri ishoda.

Izvedeno je u jeziku PHP 8.5 i okviru Laravel 13.25, kao skup konzolnih naredbi
nad bazom SQLite, bez mrežnog sloja i bez korisničkog sučelja.

Cijeli dokazni sadržaj visi o jednoj razlici u izvedbi. Zapis namjere piše se
ili prije predaje, u istoj transakciji kao i sam dokument, ili tek nakon što
odgovor stigne. Obje izvedbe stoje u istom kodu i biraju se prekidačem u
naredbi. Iz te jedne razlike slijede sva tri kontrasta koja
odjeljak~\ref{sec:scenariji} pokazuje.

% ------------------------------------------------------------
\section{Model stanja i ključne komponente}
\label{sec:komponente}
% Izvodi se IZ 6.3, ne izmišlja ovdje.
% Dva listinga: (1) tablica prijelaza + uvjet roka, (2) tumačenje S008.
% Enum stanja ide kao TABLICA, da hrvatski opis stoji uz oznaku iz koda.

Podjela na entitete iz odjeljka~\ref{sec:dizajn} izvedena je doslovno. Dokument
i njegova predaja su dva entiteta, svaki u svojoj tablici, a transakcija koja
bi ih obuhvatila ne postoji\cite{helland2007lifebeyond}. Treća tablica čuva
evidenciju. Tablica~\ref{tab:stanja} daje stanja automata predaje uz oznake
koje nose u kodu.

\begin{table}[htbp]
\centering
\caption[Stanja automata predaje]{Stanja automata predaje. Oznake su iz
programskog koda, opisi iz jezika rada.}
\label{tab:stanja}
\small
\begin{tabular}{@{}llc@{}}
\toprule
Oznaka u kodu & Značenje & Ishod poznat \\
\midrule
\texttt{prepared}            & namjera zapisana, predaja nije pokušana & da \\
\texttt{sent-unconfirmed}    & poslano, potvrda nije stigla            & \textbf{ne} \\
\texttt{confirmed}           & pristupna točka je potvrdila primitak   & da \\
\texttt{correction-rejected} & propisani ispravak odbijen šifrom S008  & da \\
\texttt{deadline-expired}    & rok istekao prije potvrde               & da \\
\bottomrule
\end{tabular}
\end{table}

Drugi redak tablice je ono zbog čega automat postoji. Stanje koje propis ne
imenuje ovdje je ravnopravno stanje, uz obrazloženje iz
odjeljka~\ref{sec:dizajn}\cite{schneider1990statemachine}.

\begin{figure}[htbp]
\centering
\begin{tikzpicture}[
  node distance=1.0cm and 1.5cm,
  st/.style={draw, rounded corners=3pt, align=center, font=\scriptsize,
             minimum width=2.1cm, minimum height=0.9cm},
  konac/.style={st, double, double distance=1pt},
  pr/.style={-{Latex[length=1.8mm]}, thick},
  ozn/.style={font=\scriptsize, align=center, inner sep=1.5pt}
]
  \node[st] (pri) {prepared};
  \node[st, right=2.1cm of pri] (nep) {sent-\\unconfirmed};
  \node[konac, right=2.3cm of nep] (pot) {confirmed};
  \node[konac, above=1.0cm of pot] (odb) {correction-\\rejected};
  \node[konac, below=1.0cm of pot] (rok) {deadline-\\expired};

  \draw[pr] (pri) -- node[ozn, above] {sent} (nep);
  \draw[pr] (nep) -- node[ozn, above] {confirmation} (pot);
  \draw[pr] (nep) -- node[ozn, above left, pos=0.35] {rejection\\from S008} (odb);
  \draw[pr] (nep) -- node[ozn, below left, pos=0.35] {deadline\\elapsed} (rok);
  \draw[pr] (nep) to[loop above] node[ozn, above] {sent} (nep);
\end{tikzpicture}
\caption[Automat predaje dokumenta]{Automat predaje. Dvostruki obrub označava
stanje iz kojeg nema prijelaza. Izostanak odgovora nije prijelaz nego petlja:
dokument koji je poslan i nije potvrđen ostaje poslan i nepotvrđen.}
\label{fig:automat}
\end{figure}

Rok iz članka 49.\ stavka 1.\ u automat ne ulazi kao stanje nego kao uvjet nad
prijelazom, iz razloga koji daje odjeljak~\ref{sec:automat}.
Listing~\ref{lst:automat} pokazuje taj uvjet.

\begin{lstlisting}[language=PHP, caption={[Tablica prijelaza i uvjet
roka]Tablica prijelaza i uvjet roka. Prijelaz pokreće poticaj, a rok samo
odlučuje je li poticaj dopušten.}, label={lst:automat}]
private const TRANSITIONS = [
    State::PREPARED->value => [
        Trigger::SENT->value => State::SENT_UNCONFIRMED,
    ],
    State::SENT_UNCONFIRMED->value => [
        Trigger::SENT->value                   => State::SENT_UNCONFIRMED,
        Trigger::CONFIRMATION->value           => State::CONFIRMED,
        Trigger::CONFIRMATION_FROM_S008->value => State::CONFIRMED,
        Trigger::REJECTION_FROM_S008->value    => State::CORRECTION_REJECTED,
        Trigger::DEADLINE_ELAPSED->value       => State::DEADLINE_EXPIRED,
    ],
];

// Rok kao uvjet nad prijelazom, ne kao stanje.
if ($trigger === Trigger::SENT && $expired) {
    throw new DomainException(
        'Ponovno slanje nije dopusteno: rok iz cl. 49. st. 1. je istrosen.'
    );
}
\end{lstlisting}

Nedopušten prijelaz baca iznimku i ne mijenja stanje. Sustav koji nedopušten
prijelaz tiho preskoči prestaje biti dokaz o vlastitom stanju.

Predaja se u odlazni pretinac zapisuje prije poziva, u istoj transakciji u
kojoj nastaje dokument, uz ključ koji određuje
pošiljatelj\cite{saltzer1984endtoend}. Zapisana namjera nosi i drugi posao,
teži od prvog. Listing~\ref{lst:s008} pokazuje mjesto na kojem se šifra S008
tumači.

\begin{lstlisting}[language=PHP, caption={[Tumačenje šifre S008 uz zapisanu
namjeru]Tumačenje šifre S008. Ista šifra vodi u dva suprotna prijelaza, a
razlučuje ih samo namjera zapisana prije slanja.}, label={lst:s008}]
return match ($intent) {
    Intent::RETRY => Interpretation::unambiguous(
        Trigger::CONFIRMATION_FROM_S008,
        'S008 uz zapisan ponovni pokusaj dokazuje da je raniji pokusaj prosao',
    ),

    Intent::CORRECTION => Interpretation::unambiguous(
        Trigger::REJECTION_FROM_S008,
        'S008 uz zapisan ispravak znaci odbijenicu propisanog tijeka ispravka',
    ),

    null => Interpretation::ambiguous(
        'S008 bez zapisane namjere: ista sifra znaci i potvrdu i odbijenicu',
    ),
};
\end{lstlisting}

Zadnji slučaj u listingu je onaj koji se ne smije prešutjeti. Kada zapisa
nema, funkcija ne pogađa nego vraća dvoznačan ishod, a automat tada nema
prijelaz na raspolaganju. Rupa ostaje vidljiva umjesto da se popuni
pretpostavkom.

U evidenciju ulazi svaki prijelaz, ali i svaki trenutak u kojem tumačenje nije
bilo jednoznačno. To drugo je ono zbog čega tablica postoji.

% ------------------------------------------------------------
\section{Scenariji i ishodi}
\label{sec:scenariji}
% ⚠ SVI ISPISI SU STVARNI, uzeti 2026-08-13. Sat prototipa je fiksiran
%   na 2026-09-01 09:00:00, pa su ponovljivi. Ako se kod mijenja,
%   ISPISE UZETI PONOVO.
%
% ★ SVAKI scenarij završava rečenicom koja imenuje ĆELIJU matrice.

Tri scenarija pokrivaju po jedan način otkazivanja iz
odjeljka~\ref{sec:failure-modes}. Svaki se pokreće u obje izvedbe, sa zapisom
namjere i bez njega. Sat je pritom zamrznut na jedan trenutak, pa su ispisi
ponovljivi.

\begin{table}[htbp]
\centering
\caption[Scenariji i završna stanja]{Scenariji i završna stanja. Razlika među
stupcima nastaje samo iz trenutka u kojem se zapisuje namjera.}
\label{tab:scenariji}
\small
\begin{tabular}{@{}llll@{}}
\toprule
& Odgovor & Sa zapisom namjere & Bez zapisa namjere \\
\midrule
B1 & izostaje       & \texttt{sent-unconfirmed}    & nema zapisa \\
E1 & izostaje       & predaja bez potvrde: 1       & predaja bez potvrde: 0 \\
B2 & S008 uz ispravak & \texttt{correction-rejected} & \texttt{sent-unconfirmed} \\
\bottomrule
\end{tabular}
\end{table}

\paragraph{B1: poziv istekne bez odgovora.}
Sa zapisom namjere dokument završava u stanju \texttt{sent-unconfirmed}, uz
zabilježen dan nastupa nemogućnosti i zakazan sljedeći pokušaj. Bez zapisa u
bazi ne postoji nijedan redak o predaji, iako je račun izdan.

Scenarij pokazuje i zašto raspored ponavljanja mora gledati rok.
Tablica~\ref{tab:raspored} stavlja dva rasporeda jedan pokraj drugog.

\begin{table}[htbp]
\centering
\caption[Raspored ponavljanja prema roku i prema broju pokušaja]{Raspored
ponavljanja uz rok koji istječe 8.~rujna u 23:59:59. Lijevi stupac gleda
preostali rok, desni samo broj dotadašnjih pokušaja.}
\label{tab:raspored}
\small
\begin{tabular}{@{}ccc@{}}
\toprule
Pokušaj & Svjestan roka & Eksponencijalni odmak \\
\midrule
1 & 2026-09-05 04:29:59 & 2026-09-01 09:01:00 \\
2 & 2026-09-07 02:14:59 & 2026-09-01 09:03:00 \\
3 & 2026-09-08 01:07:29 & 2026-09-01 09:07:00 \\
4 & 2026-09-08 12:33:44 & 2026-09-01 09:15:00 \\
5 & 2026-09-08 18:16:51 & 2026-09-01 09:31:00 \\
\bottomrule
\end{tabular}
\end{table}

Nalaz je precizniji od očekivanog. Eksponencijalni odmak rok ne probija rano
nego ga uopće ne dodiruje. Prvih pet pokušaja potroši u trideset i jednoj
minuti, a rok prekorači tek na četrnaestom. Mana obrasca ovdje nije da kasni.
Mana je da mu raspored s rokom nema nikakve veze, pa se isti obrazac uz
drukčiju osnovicu ponaša posve drukčije, a propis se u tom izboru ne
pojavljuje. Raspored svjestan roka svaki pokušaj stavlja na polovicu
preostalog roka i zato ga po konstrukciji ne može prekoračiti. Brookerovo
slučajno raspršivanje razmaka isključeno je, da ispis ostane
ponovljiv\cite{brookertimeouts}.

Scenarij popunjava ćelije B1 uz automat s imenovanim stanjem i B1 uz
ponavljanje svjesno roka.

\paragraph{E1: predaja bez potvrde.}
Pitanje na koje scenarij odgovara nije je li dokument stigao. Pitanje je može
li poslovni sustav znati da ga je pokušao predati. Sa zapisom namjere odgovor
postoji i provjerljiv je bez pristupa pristupnoj točki. Bez zapisa račun
postoji, a trag predaje ne postoji.

Scenarij popunjava ćeliju E1 uz odlazni pretinac i zapis namjere.

\paragraph{B2: ista šifra za dva ishoda.}
Scenarij ima dva dijela. U prvom se poruka ponavlja nakon isteka vremena i
dobiva S008. U drugom se šalje ispravak pod istim brojem računa i dobiva istu
šifru. Listing~\ref{lst:ispis} prepisuje završno stanje baze.

\begin{lstlisting}[caption={[Završno stanje baze nakon scenarija
B2]Završno stanje nakon scenarija B2, izlaz naredbe \texttt{php artisan
state}. Ista šifra, dva različita završna stanja.}, label={lst:ispis}]
PREDAJE
+----+------------+------------+---------------------+------+---------+
| id | racun      | namjera    | stanje              | pok. | odgovor |
+----+------------+------------+---------------------+------+---------+
| 1  | R-2026-003 | retry      | confirmed           | 2    | s008    |
| 2  | R-2026-003 | correction | correction-rejected | 1    | s008    |
+----+------------+------------+---------------------+------+---------+

EVIDENCIJA
+---------+-----+------------------------+------------------+-------------+
| predaja | br. | dogadaj                | u                | jednoznacno |
+---------+-----+------------------------+------------------+-------------+
| 1       | 1   | zapisana namjera       | prepared         | da          |
| 1       | 2   | sent                   | sent-unconfirmed | da          |
| 1       | 3   | tumacenje odgovora     | sent-unconfirmed | NE          |
| 1       | 4   | sent                   | sent-unconfirmed | da          |
| 1       | 5   | confirmation-from-s008 | confirmed        | da          |
+---------+-----+------------------------+------------------+-------------+
\end{lstlisting}

Treći redak evidencije je onaj koji nosi tezu poglavlja. Sustav je u tom
trenutku znao da ne zna, i to je zapisao. Kasniji odgovor S008 taj trenutak
razrješava, ali ga ne briše.

Bez zapisa namjere oba dijela scenarija završavaju u istom stanju,
\texttt{sent-unconfirmed}. Dvije suprotne namjere, ista šifra i isti ishod.
To je slučaj B2 u izvedbi.

Scenarij popunjava ćeliju B2 uz odlazni pretinac i zapis namjere.

% ------------------------------------------------------------
% TVRDNJE KAO TESTOVI — veza s metodom iz pogl. 8.
% CoFI [chen2020cofi] i Filibuster [meiklejohn2021filibuster] su citirani
% u 8.1 kao opravdanje metode; ovdje se citira samo OBLIK tvrdnje.

Svaka tvrdnja iz ovog odjeljka zapisana je i kao izvršni test, u obliku koji
Filibuster koristi za uvjetne tvrdnje\cite{meiklejohn2021filibuster}. Test ne
provjerava da je poziv uspio nego da je završno stanje ono koje propis i
obrazac zajedno nalažu. Prototip nosi dvadeset i tri takve tvrdnje, a
poglavlje~\ref{ch:evaluacija} poziva se na njih kada ćeliju označava kao
pokazanu.

